\section{The decision estimand}
\label{sec:decision}
%==========================================================================

\emph{Section~\ref{sec:designsec} produces an identified $\theta$. This section asks what it
decides, and the answer is the negative result the paper ends on.}


Confidence intervals on $\theta$ are the wrong output. The clinical question is the
\emph{sign} of $\mu = \theta - \thc$, and its two errors are not interchangeable: treating a
patient who does not need it costs an infusion and a venous line, and failing to treat one
who does costs a decompensation. Figure~\ref{fig:D3} is the operating characteristic of
that decision, and it carries the negative result the paper ends on. Everything above says
the therapy channel makes $\theta$ identifiable; it does not follow that it makes the sign
of the margin decidable, and it does not.

\begin{figure}[htbp]
\centering
\includegraphics[width=\textwidth]{fig_D3_decision.pdf}
\caption{The probability of getting the sign of $\mu$ wrong, by simulation, as a patient of
fixed biology is swept across the decision boundary by the illness amplitude. Left of the
dotted line the error is a patient treated who did not need it; right of it, a patient left
untreated who did. Each panel varies one thing a protocol controls.}
\label{fig:D3}
\end{figure}

\begin{table}[htbp]
\centering
\footnotesize
\caption{\label{tab:decision}The decision characteristic, by simulation, on three designs. Twelve seeds at four hundred replicates each; the $\pm$ is the spread across seeds, which twelve seeds fix only to about a fifth of its own value, so it is a scale and not two significant figures. Both error columns are taken at $|\mu| = 339$ either side of the boundary.}
\begin{tabular}{@{}lrrr@{}}
\toprule
\textbf{design} & \textbf{treat unneeded} & \textbf{fail to treat} & \textbf{cohort $N$} \\
\midrule
fixed operating point, $n=9$ & $6.9\% \pm 1.7$ & $44.0\% \pm 3.2$ & $4900$ \\
monthly across dose cycles, $n=25$ & $19.1\% \pm 2.1$ & $41.6\% \pm 2.5$ & $2200$ \\
\quad the same, increment known only to $11.6\%$ & $19.1\% \pm 1.7$ & $42.3\% \pm 2.9$ & n/a \\
\bottomrule
\end{tabular}
% twelve seeds at four hundred replicates each, the spread across seeds; the two error columns are at |mu| = 339 either side of the boundary
\end{table}

Table~\ref{tab:decision} puts the same decision to three designs. Near the boundary the
therapy-channel design errs on about two patients in five, against four in nine for a
design fixed at the operating point. The cohort needed to place the fraction of
patients with $\mu<0$ within ten points falls from about $\num{cohort_fixed}$ to about $\num{cohort_monthly}$: a
real gain, and still a number no cohort of this disease reaches, since the trial that
established the therapy enrolled sixteen~\cite{trial_phase12} and
Section~\ref{sec:howmany} sets $\num{cohort_monthly}$ against the size of the diagnosed population
rather than against trial logistics. % C16
The dose-precision requirement does not bind on this estimand: degrading the increment
from exact to the $\num{kappa_monthly_pct}\%$ ($\num{kappa_monthly_pct_lo}$ to $\num{kappa_monthly_pct_hi}$) the identifiability gate asks for moves the
sign error by less than its seed-to-seed spread. And the reason the decision does not
improve is not identifiability at all, but that the sign of a quantity near its own zero
is hard for reasons independent of what the record contains. These numbers do not
contradict Table~\ref{tab:designs}, which reports a posterior standard deviation under
the $k_e$ prior where this table reports a prior-free nonlinear least squares; the gap
between them is the shape of the result. The therapy channel identifies $\theta$ given a
defensible prior on the alternative channel, and does not make the decision reliable under
any reading.

\paragraph{The asymmetry reverses under crisis protocol.} The estimator regresses the
\emph{recorded} load on the marker. Write the recorded load as
$L_{\mathrm{base}} + \rho\,(s_{\mathrm{true}} - L_{\mathrm{base}})$, so $\rho = 1$ is a
record that tracks the load to its recording error and $\rho \le 0$ is crisis protocol,
where prescribed intake is cut while catabolic muscle breakdown supplies the load. On the
monthly design $\rho = 1$ errs at $17.5\%$ and $41.8\%$ with a median $\hat\theta/\theta$
of $1.03$; at $\rho = 0.5$ the two errors become $99.0\%$ and $0.0\%$ at $0.90$; at
$\rho = 0$, intake pinned at baseline, $100\%$ and $0.0\%$ at $0.086$; at $\rho = -0.25$,
protein stopped, $100\%$ and $0.0\%$ at $0.009$. The mechanism is in the last of those
numbers: when recorded intake falls while the marker climbs, least squares has nowhere to
put the marker except lost capacity, and $\hat\theta$ collapses by two orders of
magnitude. ``The estimator fails towards declaring patients safe'' holds only while the
record tracks the load; under the protocol a decompensation actually triggers, the same
estimator declares every patient unsafe. Neither half belongs in an abstract on its own
and neither is a bug to be fixed by weighting: they are one estimator meeting two
data-generating processes, and which one a record came from is a fact about the chart
rather than about the patient. Nor does a regret-equalised threshold repair it. At equal
costs the managed regime wants the boundary at $+0.07\,\theta$ and the crisis regime wants
it at $-0.39\,\theta$; at a cost ratio of ten to one or worse every regime wants every
patient escalated, which is the decision-theoretic way of saying the test has stopped
contributing.

These error rates are simulated rather than propagated, and the difference matters. The
design's near-null direction is very nearly the direction along which $\mu$ is monotone
and of one sign, so the sampling distribution of $\hat\mu$ is skewed and does not cross
zero anything like as often as its standard deviation suggests: a delta-method calculation
reporting $\Phi(-|\mu|/\operatorname{sd}(\hat\mu))$ would throw away the one favourable
result this paper has. That is the margin theorem of the supplement's fibre calculation
(\S S1) showing up in the estimator, and both numbers are reported by the code so the gap
stays visible. A linearisation understates what the record contains here and overstates it
in the design matrix of Section~\ref{sec:rank}, which is the lesson: a diagnostic is a
model too.


\vspace{0.5em}
\hrule
\vspace{1em}

%==========================================================================
